\section{Related Work}
\label{sec:related}

We survey three lines of work that collectively motivate but do not provide
structural efficiency measurement for code generation alignment.

\subsection{Execution-Feedback Reinforcement Learning for Code Generation}

A growing body of work demonstrates that execution feedback improves code
generation beyond supervised fine-tuning. CodeRL \cite{Le2022CodeRL} introduces
actor-critic training with unit test feedback, establishing that execution
outcome signals can guide policy improvement. PPOCoder \cite{Shojaee2023Execution}
extends PPO-based RL to code synthesis, showing gains on HumanEval and MBPP.
T\"{U}LU~3 \cite{Lambert2024TULU} demonstrates that GRPO with binary pass/fail
rewards on code problems yields competitive pass@1 improvements in a
multi-task alignment setting.

\textbf{The limitation of this line:} all comparisons are made on pass@1 outcomes.
None of these works measures \emph{how} the policy moves structurally relative to
the base model. It is unknown whether the observed improvements reflect
genuinely richer structural learning or more decisive surface-level pattern
matching. Our framework provides the diagnostic to answer this question.

\subsection{DPO and Preference-Based Alignment}

Direct Preference Optimization \cite{Rafailov2023Direct} provides an alternative
post-training objective that avoids online rollouts by framing RLHF as a
supervised learning problem over preference pairs. DPO has been applied to
code generation \cite{Guo2024DeepSeek,Luo2023Execution}, typically using
execution-oracle labeling (passing vs.\ failing solutions as preferred/rejected).
GRPO \cite{Shao2024DeepSeekMath} reformulates PPO's policy gradient without a critic
network, using group-relative rewards, and has become the dominant execution-RL
method for code alignment following DeepSeek-R1 \cite{DeepSeekAI2025DeepSeek}.

\textbf{The limitation of this line:} DPO and GRPO are compared on benchmark outcomes,
not on structural properties of the policy change they induce. KL divergence
appears in both frameworks as a constraint (DPO's implicit KL regularization;
GRPO's KL penalty term $\beta$), but is used only as a training-time constraint,
never as a diagnostic normalizer for structural movement analysis. We repurpose
KL divergence from a training constraint into a measurement tool.

\subsection{AST-Based Code Analysis and Similarity Metrics}

The program analysis community has developed rich tools for AST-level code
comparison. The FA-AST taxonomy \cite{Wang2020Detecting} classifies Python AST nodes
into control-flow (\texttt{If}, \texttt{For}, \texttt{While}, \texttt{Try}, \texttt{With}), data-flow (\texttt{Assign}, \texttt{Call}, \texttt{Return},
\texttt{FunctionDef}), and surface categories, providing a principled semantic
classification of structural program elements. The ZSS algorithm \cite{Zhang1989Simple}
computes minimum-cost tree edit distance and has been validated as
correlating with established code similarity metrics \cite{Ding2024Revisiting}.
Code clone detection \cite{Svajlenko2014Towards} and program
repair \cite{Hua2018SketchFix} literatures use AST edit distances extensively.

\textbf{The limitation of this line:} AST analysis in program analysis is applied to
pairs of programs, not to the \emph{policy movement} of a trained model relative to
its base. The composition of FA-AST classification, ZSS edit distance,
and KL-matched checkpoint comparison --- which enables structural efficiency
measurement --- has not been proposed or validated.

\subsection{Policy Analysis and Interpretability in RL Fine-Tuning}

Recent work on understanding RL fine-tuning \cite{Schulman2017Proximal,Zheng2023Secrets}
focuses on training stability, reward hacking, and KL divergence management.
Mechanistic interpretability work \cite{Anthropic2023Interpretability,Elhage2022Circuits} analyzes
model internals (attention heads, circuits) but does not specifically measure
the structural quality of code-level policy movement as a function of alignment
method.

\textbf{Our position:} We bring program analysis tools (FA-AST, ZSS) into the alignment
evaluation setting, providing a structural lens on policy movement that complements
both activation-level interpretability and benchmark-based evaluation. The result
is a framework that any team can apply to their checkpoint archive to diagnose
whether their alignment method produces structurally rich policy change.
